% Ad Familiares 10.18 (ad-familiares-10-18) --- English translation
% Translated by Claude (Anthropic) from the Latin (Perseus).
% Style guide: STYLE.md.

\ciceroLetterOpener{PLANCVS CICERONI}

\ciceroSection{1} What I had in mind when Laevus and Nerva left me, you could learn both from the letter I gave them and from the men themselves, who shared in all my doings and counsels. It has fallen to me to do what is apt to fall to a man of self-respect who is eager to give satisfaction to the commonwealth and to all loyal men: I chose the path of greater risk, so long as I could give a good account of myself, rather than the safe one, which could give the carpers something to fasten on.

\ciceroSection{2} And so, after the legates' departure, when two letters in unbroken succession came from Lepidus asking me to come, and from Laterensis pleading and adjuring me much more urgently still --- shrinking from no other thing than that very thing which is the source of my own fear too, the inconstancy and unreliability of Lepidus's army --- I judged I could not hesitate to come to their aid and put myself in the path of the common danger. For I saw that, although the more cautious course was to wait at the Is\`ere until Brutus brought his army across, and then to march out to meet the enemy with a colleague of one mind with me, with an army at peace within itself and right-thinking about the state, the way soldiers ought to be --- still, if Lepidus, with his right thinking, suffered any harm, the whole would have been laid to my account, either as obstinacy or as fear, on the ground that I had failed either to lift up a man whose feelings had been wounded though he was firm with the state, or had myself drawn back from sharing in a struggle so necessary.

\ciceroSection{3} So I chose rather to run the risk --- to see whether by my presence I could both safeguard Lepidus and make the army better --- than to look over-cautious. Certainly no man without his own affairs constraining him, I think, has ever been more anxious. For what would have given no occasion for hesitation, if Lepidus's army were out of the picture, brings now a great anxiety and great hazard. For if it had been my luck to come up first against Antony, by Hercules he would not have stood his ground for an hour --- such is my confidence both in myself and in the despicable state of his forces and the muleteer Ventidius's camp; but I cannot help shuddering, if some wound is festering beneath the skin which can do its damage before it can be either known or treated. Yet for certain, unless I hold myself in one place, Lepidus himself would run great danger, and so would that part of the army which thinks rightly about the state. Even the desperate enemy would have gained a great accession of strength, had they been able to draw off any of Lepidus's troops. If my arrival checks all this, I shall give thanks to fortune, and to my own steadiness, which has roused me to make this trial.

\ciceroSection{4} So on the 18th of May I broke camp from the Is\`ere; the bridge, however, which I had built across the Is\`ere I have left in place, with two forts set at its heads, and have put strong garrisons there, so that, when Brutus comes with his army, the crossing may be ready without delay. As for myself, I hope to join Lepidus's forces within the next eight days from the day on which I am dispatching this letter.
